\documentclass[UTF8,fontset=fandol,11pt]{ctexart}
\usepackage[a4paper,margin=25mm]{geometry}
\usepackage{amsmath,amssymb,amsthm,mathtools}
\usepackage{xcolor}
\usepackage{booktabs,array,enumitem,fancyhdr,hyperref}
\hypersetup{colorlinks=true,linkcolor=blue!50!black,urlcolor=blue!50!black}
\setlength{\emergencystretch}{2em}
\setlength{\parskip}{0.3em}
\linespread{1.16}
\allowdisplaybreaks
\newtheorem{theorem}{定理}[section]
\newtheorem{proposition}[theorem]{命题}
\newtheorem{lemma}[theorem]{引理}
\theoremstyle{definition}
\newtheorem{definition}[theorem]{定义}
\theoremstyle{remark}
\newtheorem{remark}[theorem]{说明}
\newcommand{\dd}{\,\mathrm d}
\newcommand{\ii}{\mathrm i}
\newcommand{\Rea}{\operatorname{Re}}
\newcommand{\Ima}{\operatorname{Im}}
\newcommand{\Li}{\operatorname{Li}}
\newcommand{\supp}{\operatorname{supp}}
\pagestyle{fancy}
\fancyhf{}
\fancyhead[L]{\small 学习笔记}
\fancyhead[R]{\small 初稿}
\fancyfoot[C]{\thepage}
\setlength{\headheight}{15pt}
\author{}
\date{2026年10月22日初稿}
\makeatletter
\renewcommand{\maketitle}{%
  \thispagestyle{plain}%
  \begin{center}
    {\LARGE\bfseries \@title\par}
    \vspace{0.6em}
    {\small \@date\par}
  \end{center}
  \vspace{0.6em}}
\makeatother
\newcommand{\draftnotice}{\noindent{\small
本文为供逐篇审阅的学习笔记初稿，不是新增研究成果。
文中区分直接证明与引用结果；正文不依赖本文。
本轮已作公式、条件及排版初步核查，尚待作者逐句审定。}\par\medskip}

\title{从矩形求和到平滑求和}
\begin{document}
\maketitle
\draftnotice
\section{改变的不是系数，而是截断的方式}
给定系数 $a_n$，普通部分和 $\sum_{n\le x}a_n$ 在 $n=x$ 处突然停止。
另一种选择是让接近上限的项逐渐减小：
\[
 A_1(x)=\sum_{n\le x}a_n(1-n/x).
\]
这不是普通部分和的另一种记号，而是新的加权对象。
本篇研究权重如何决定积分核，以及多一个分母因子究竟改善了什么。

\section{从权重算出核}
对非负整数 $m$，令
\[
 W_m(u)=
 \begin{cases}(1-u)^m/m!,&0<u<1,\\0,&u>1.\end{cases}
\]
当 $m\ge1$ 时再置 $W_m(1)=0$。对 $\Rea s>0$，
\[
 \int_0^1u^{s-1}\frac{(1-u)^m}{m!}\dd u
 =\frac1{s(s+1)\cdots(s+m)}.
\]
证明不必预先使用 Beta 函数。$m=0$ 时积分为 $1/s$；
$m\ge1$ 时分部积分，端点项在 $0$ 处因 $\Rea s>0$ 消失，在 $1$ 处因权重消失，并得
\[
 I_m(s)=\frac1s I_{m-1}(s+1).
\]
递推即可得到上式。

对 $m\ge1$ 和 $c>0$，权重在 $u=1$ 连续，且其 Mellin 变换沿竖线绝对可积。
由 Mellin 反演并取 $u=1/y$，
\[
 \frac1{2\pi\ii}\int_{c-\ii\infty}^{c+\ii\infty}
 \frac{y^s}{s(s+1)\cdots(s+m)}\dd s
 =
 \begin{cases}
 (1-1/y)^m/m!,&y>1,\\
 0,&0<y\le1.
 \end{cases}
\]
在 $y=1$ 处取零，来自连续性，不是普通阶跃的半权规则。
$m=0$ 时不满足这里的绝对可积条件，须改用对称截断的阶跃核。

\section{对 Dirichlet 级数逐项反演}
设 $F(s)=\sum_{n\ge1}a_nn^{-s}$，并选择 $c>0$ 使
$\sum |a_n|n^{-c}<\infty$。对 $m\ge1$，
\[
 A_m(x):=\sum_{n\le x}a_n\frac{(1-n/x)^m}{m!}
 =\frac1{2\pi\ii}\int_{c-\ii\infty}^{c+\ii\infty}
   F(s)\frac{x^s}{s(s+1)\cdots(s+m)}\dd s.
\]
换序依据是
\[
 x^c\sum_{n\ge1}|a_n|n^{-c}
 \int_{\mathbb R}
 \frac{\dd t}{|(c+\ii t)(c+1+\ii t)\cdots(c+m+\ii t)|}<\infty.
\]
分母沿竖线为 $O(|t|^{m+1})$，而有限高度范围内没有分母零点。
所以这里对级数与积分的交换有整体绝对可积依据，不依靠形式逐项操作。

\section{一次平滑等于一次积分}
令 $A(t)=\sum_{n\le t}a_n$，并在 $0\le t<1$ 取零。
由于任意有限区间上只有有限项，
\[
 C_1(x):=xA_1(x)=\sum_{n\le x}a_n(x-n)
       =\int_0^xA(t)\dd t.
\]
每个跳跃被积分成连续的折线段。在非整数开区间内，$C_1'(x)=A(x)$；
在有非零系数的整数处，两侧斜率一般不同，不能称处处可微。
重复积分则产生更高次权重：
\[
 C_m(x):=\frac1{m!}\sum_{n\le x}a_n(x-n)^m=x^mA_m(x),\qquad
 C_m'(x)=C_{m-1}(x)\quad(m\ge2).
\]

\section{“改善收敛”需要说明比较对象}
在原积分线 $\Rea s=c$ 上，$F$ 有界。对一次平滑的尾积分，有
\[
 \left|\frac1{2\pi\ii}\int_{c+\ii T}^{c+\ii\infty}
 F(s)\frac{x^s}{s(s+1)}\dd s\right|
 \le\frac{x^c}{2\pi}
 \left(\sum |a_n|n^{-c}\right)\int_T^\infty\frac{\dd t}{t^2}
 =O_{c,x,a}(T^{-1}).
\]
下半条尾线同理。这说明沿原线的截断误差容易控制。
但移线后 $F(s)$ 可能增长或有极点，仍须重新估计，不能仅凭多一个 $s+1$ 就宣布所有边积分消失。

平滑还有代价：$C_1$ 给的是累计信息，不是原来的 $A$。
若 $a_n\ge0$，则 $A$ 单调，并有
\[
 \frac{C_1(x)-C_1(x-h)}h\le A(x)
 \le\frac{C_1(x+h)-C_1(x)}h\qquad(0<h<x).
\]
若 $C_1$ 的近似误差为 $E$，差商会把误差放大到约 $E/h$。
选 $h$ 时须同时控制主项变化和这一误差；恢复不平滑结论并非免费。

\section*{参考依据与后续核验}
Mellin 反演作为标准定理引用；本篇给出了核计算、换序及一次平滑的全部基本步骤。
本书的 Perron 表示可作为背景。下一步可用 $a_n=1$ 与 $a_n=\Lambda(n)$ 分别做数值对照，
但本文不以实验替代上述估计。
\end{document}
